\documentclass[10pt,a4paper]{amsart}
\usepackage[margin=19mm]{geometry}
\usepackage{amsmath,amssymb,mathtools}
\usepackage[scr=rsfs,cal=euler]{mathalfa}
\usepackage{xfrac}
\usepackage[unicode,hidelinks,bookmarks=false]{hyperref}
\newcommand{\ie}{i.e.\ }
\newcommand{\cf}{cf.\ }
\newcommand{\resp}{respectively\ }
\newcommand{\Subsection}{\subsection}
\providecommand{\nobreakdash}{\nobreak\hbox{-}}
\setlength{\emergencystretch}{2em}
\multlinegap0pt
\usepackage{bm}
\usepackage{bbm}
\usepackage{dsfont}
\usepackage{xspace}
\usepackage{cancel}
\usepackage{mathtools}
\usepackage{amsthm}
\makeatletter
\@ifundefined{lemma}{\newtheorem{lemma}{Lemma}}{}
\@ifundefined{theorem}{\newtheorem{theorem}{Theorem}}{}
\makeatother
\def\sE{{\mathscr{E}}}
\def\supp{\hbox{supp}\, }
\title[Asymptotic behavior of solutions to a space fractional diffu]{Asymptotic behavior of solutions to a space fractional diffusion equation}
\author{Barbara {\L}upi\'nska, Piotr Rybka}
\date{Source: arXiv:2511.06030v1 (CC BY 4.0)}
\begin{document}
\maketitle

\noindent\emph{Source and licence.} Asymptotic behavior of solutions to a space fractional diffusion equation. Version arXiv:2511.06030v1 (CC BY 4.0), distributed under the Creative Commons Attribution 4.0 International licence, as recorded in the arXiv licence metadata for this exact version and shown on the abstract page https://arxiv.org/abs/2511.06030v1. Full rights notice: licenses/arxiv-2511.06030-CC-BY-4.0.txt.\par\smallskip

\noindent\textit{Editorial scope.} Counted unit: the Lemma whose source label is \texttt{l5} in the article (source0.tex lines 349--355, source label \texttt{l5}), delivered together with the article's own complete original proof of it (source0.tex lines 356--452, one \texttt{proof} environment of 97 lines). No mathematics is written, repaired or re-derived: statement and proof are the article's own text line for line. Required context retained uncounted: df-e (lines 167--170); l4 (lines 217--234); poch (lines 220--229); rn3.5 (lines 304--306); rn3.6 (lines 307--309); t1 (lines 321--342); df-m (lines 323--325); df-v (lines 344--347). Every definition, display and label of those lines is retained unchanged. Omitted from this package and not redistributed: 1--166, 171--216, 230--303, 310--320, 326--343, 348--348, 453--555. Nothing inside a retained excerpt is omitted or altered: every retained body is delivered exactly as the article's lines are written, so no omission receipt of any kind accompanies this package. Citations: the retained excerpts carry no citation command at all, so no bibliography entry of the article is transcribed and no reference is redistributed. Reference adaptation: none was needed and none was made; every reference inside the delivered unit resolves inside it, so no unresolved reference or citation is produced. Printed slips kept literal: no printed word, symbol, hypothesis or number of the counted statement or of its proof was altered. Layout and class: the article's own class is \texttt{article} with packages including graphicx, amsmath, amsthm, comment, amsfonts, mathrsfs, authblk, anysize, color, xcolor; the wrapper is the lane's pinned \texttt{amsart} editorial wrapper at 10pt on A4, which restates the theorem-like environments the retained text needs and adds the article's own macro definitions that the retained text needs (\texttt{\textbackslash sE}, \texttt{\textbackslash supp}), with extra wrapper packages (bm, bbm, dsfont, xspace, cancel, mathtools). The delivered numbering is the unit's own. Assets: no retained span contains a figure environment, a graphics-inclusion command, a PDF-page inclusion, a table or a TikZ picture; no image file is copied into this package, the package declares no asset dependency and no image file is redistributed with it. Licence: version arXiv:2511.06030v1 of this article is distributed under the Creative Commons Attribution 4.0 International licence, as recorded in the arXiv licence metadata for this exact version and shown on the abstract page https://arxiv.org/abs/2511.06030v1. A full rights notice is delivered at licenses/arxiv-2511.06030-CC-BY-4.0.txt. Characters: every retained excerpt is pure ASCII, so the delivered engine, XeLaTeX with its default Latin Modern text font, reports no missing character.
\begin{equation}\label{df-e}
\sE_t(x)=%\sE(x,t)=
a_0t^{-\frac{1}{1+\alpha}}\Phi\left(\frac{x}{t^{\frac{1}{1+\alpha}}}\right).
\end{equation}

\begin{lemma}\label{l4}
For all $x>0$, the derivative of the profile function $\Phi$ can be %defined in~\eqref{fi} satisfies the ng 
estimated as follows,
\begin{equation}\label{poch}
|\Phi'(x)|\leq 
\left\{
\begin{array}{cc}
 \frac1{\Gamma(1+\alpha)} x^\alpha    &  x\in [0, 1], \\
 %2^{\alpha+1}\Gamma(\alpha+2) x^{-2-\alpha}
 \frac{C_1(\alpha)}{x}&  x>1,
\end{array}
\right.
\end{equation}
where 
$$
C_1(\alpha) = \alpha 2^{1+\alpha}\max\bigg\{\int_{\frac12}^1 \frac{du}{ u^{1+\alpha}(1-u)^{1-\alpha}}, 2\bigg\}.
$$
\end{lemma}

\begin{equation}\label{rn3.5}
w_1(x,t) = \int_0^\infty ( \sE_t(x-y) - \sE_t(x+y))g(y)\,dy,
\end{equation}

\begin{equation}\label{rn3.6}
w_2(x,t)  = \int_0^\infty ( \sE_t(x-y) + \sE_t(x+y))g(y)\,dy.
\end{equation}

\begin{theorem}\label{t1}
Let us suppose that  $g,\ yg(\cdot) \in L^1(0,\infty)$ have a compact support, $\theta$ is any fixed number from the interval $(0,1)$, 
\begin{equation}\label{df-m}
m=\int_0^\infty g(y)dy,
\end{equation} 
and $w_1$, (resp. $w_2$), is defined by~\eqref{rn3.5}, (resp.~\eqref{rn3.6}).
Then, \\
(a) there is $C_1>0$ such that 
$$
\|w_1(\cdot,t) \|_{L^{1/(1-\theta)}} \le C_1  t^{- (1+\theta)/(1+\alpha)} \|y g(\cdot)\|_{L^1}, \qquad t>1;
$$
(b) there is $C_2>0$ such that 
$$
\|w_2(\cdot,t)-2m\sE_t(\cdot)\|_{L^{1/(1-\theta)}} \le C_2 
 t^{- (1+\theta)/(1+\alpha)} \|y g(\cdot)\|_{L^1},\qquad t>1;
$$
(c) for $C_2(\alpha)$ defined in Lemma \ref{l5} we have
$$
\|w_2(\cdot,t)-2m\sE_t(\cdot)\|_{L^\infty}, \quad 
\|w_1(\cdot,t) \|_{L^\infty} \le \frac{C_2(\alpha)}{t^{2/(1+\alpha)}} \|y g(\cdot)\|_{L^1}, \qquad t>1.
$$
\end{theorem}

\begin{equation}\label{df-v}
v^+(x,t) = \int_0^\infty    \sE_t(x+y)g(y)\,dy,\qquad
v^-(x,t) = \int_0^\infty \sE_t(x-y) g(y)\,dy.
\end{equation}

\begin{lemma}\label{l5}
Let us suppose that $g$ satisfies the assumptions of Theorem \ref{t1}, $m$ is given by \eqref{df-m} and $v^+$, $v^-$ are defined in \eqref{df-v}. Then, 
$$
|v^\pm (x, t)  - m \sE_t(x)| \le \frac{C_2(\alpha)}{t^{1/(1+\alpha)}}\| y g(\cdot)\|_{L^1}\bigg(\frac1{t^{1/(1+\alpha)}}\chi_{[0, t^{1/(1+\alpha)}]}(x) + \frac 1x \chi_{(t^{1/(1+\alpha)}, \infty)}(x)\bigg)\qquad\hbox{for }t>1,
$$
where $C_2(\alpha) =2a_0 C_1(\alpha)$ and $C_1(\alpha)$ is defined in Lemma \ref{l4}.
\end{lemma}

\begin{proof} We will first estimate $v^+ - \sE_t$.
Let us notice that due to the definitions of $\sE_t$ and $m$, see \eqref{df-e} and \eqref{df-m}, we have
\begin{align*}
v^+(x,t) - m\sE_t(x) &= \frac{a_0}{t^{1/(1+\alpha)}}
\int_0^\infty
\bigg( \Phi \bigg(\frac{x+y}{t^{1/(1+\alpha)}}\bigg)
-  \Phi \bigg(\frac{x}{t^{1/(1+\alpha)}}\bigg)\bigg)g(y)\, dy\\
&= \frac{a_0}{t^{1/(1+\alpha)}}
\int_0^\infty \int_0^1 \frac{y}{t^{1/(1+\alpha)}}
\Phi'\bigg(\frac{x+sy}{t^{1/(1+\alpha)}}\bigg)\, ds g(y)\, dy =:I
.
\end{align*}
We consider two cases: $\frac{x}{t^{1/(1+\alpha)}}< 1$ and  $\frac{x}{t^{1/(1+\alpha)}} \ge 1$. 

When $\frac{x}{t^{1/(1+\alpha)}}< 1$ we set $y_0 = t^{1/(1+\alpha)}- x$ and $\xi=\frac{x+sy}{t^{1/(1+\alpha)}}$. We calculate
$$
I = \frac{a_0}{t^{2/(1+\alpha)}} \bigg(\int_0^{y_0} + \int_{y_0}^\infty\bigg)\int_0^1 y \Phi'\big(\xi%frac{x+sy}{t^{1/(1+\alpha)}}
\big)\, ds g(y)\, dy = I_1 + I_2.
$$
We first estimate $|I_1|$ by noticing that $|\xi|\le1$ and $x+y_0= t^{1/(1+\alpha)}$,
\begin{align*}
 |I_1| &\le \frac{a_0}{\Gamma(1+\alpha)t^{2/(1+\alpha)}} \int_0^{y_0} \int_0^1y| \xi|^\alpha
% \bigg(\frac{x+sy}{t^{1/(1+\alpha)}}\bigg)^\alpha
\, ds |g(y)|\, dy
\le \frac{a_0}{\Gamma(1+\alpha)t^{2/(1+\alpha)}} 
\int_0^{y_0} y| g(y)| \, dy,   
\end{align*}
where we applied  Lemma \ref{l4} to bound $\Phi'.$

Now, we estimate $I_2$  and again we will use Lemma \ref{l4},
\begin{align*}
 |I_2 |&\le 
 \frac{a_0 C_1(\alpha)}%\Gamma(1+\alpha)
 {t^{2/(1+\alpha)}} \int_{y_0}^\infty \bigg(
 \int_0^{s_0(y)} \bigg(\frac{x+sy}{t^{1/(1+\alpha)}}\bigg)^\alpha \, ds + \int_{s_0(y)}^1 
 \frac {t^{1/(1+\alpha)}}{x+sy}\,ds \bigg)y|g(y)|\, dy\\ &
 %= \frac{a_0 C(\alpha)}{\Gamma(1+\alpha)t^{1/(1+\alpha)}} \int_{y_0}^\infty \ln(1 + \frac yx)g(y)\, dy\\&
 \le \frac{a_0 C_1(\alpha)}{%\Gamma(1+\alpha)
 t^{2/(1+\alpha)}} \int_{y_0}^\infty y |g(y)|\, dy,
\end{align*}
where we set $s_0(y) = (t^{1/(1+\alpha)}-x)/y$ for $y\ge y_0$. %used $\ln(1+a) \le a$ for $a>0$.
Hence,
$$
|I| \le \frac{a_0 C_1(\alpha)}{%\Gamma(1+\alpha)
t^{2/(1+\alpha)}} \int_{0}^\infty y |g(y)|\, dy
$$
for $x< t^{1/(1+\alpha)}$.

Now, we take care of the second case, i.e. $\frac{x}{t^{1/(1+\alpha)}} \ge 1$. If this happens, then the argument of $\Phi'$ is always greater than 1. Hence, we proceed as in the estimate of $I_2$ above, however, $y_0=0$. Thus,
$$
|I| \le \frac{a_0 C_1(\alpha)}{%\Gamma(1+\alpha)
xt^{1/(1+\alpha)}} \int_{0}^\infty y |g(y)|\, dy.
$$
Our claim follows for $v^+$.

%These pointwise estimate lead us to the desired result for $v^+ - m\sE_t$. 

Now, we turn our attention to $v^- -m\sE_t$. We notice
$$
v^-(x,t) - m\sE_t (x)= \frac{a_0}{t^{1/(1+\alpha)}}\int_0^\infty
\bigg(\Phi\bigg( \frac{x- y}{t^{1/(1+\alpha)}}\bigg)
-\Phi\bigg( \frac{x}{t^{1/(1+\alpha)}}\bigg)\bigg)g(y)\,dy.
$$
We also see that $\Phi$ extended to negative arguments is differentiable at $0$ due to \eqref{poch}. Hence,
$$
v^-(x,t) - m\sE_t(x) = -\frac{a_0}{t^{2/(1+\alpha)}}\int_0^\infty
\int_0^1 y \Phi'\bigg(\frac{x-sy}{t^{1/(1+\alpha)}}\bigg)\,ds g(y)\, dy=:I.
$$
We must take into account the possibility that  the argument of the integrand  is negative. We recall that our datum $g$ has a compact support, we may assume that $\supp g\subset[0,R]$. Since we are interested in the time assymptotics we may consider only $t> R^{1+\alpha}.$
We will come up with different estimates for $|v^-(x,t)-m \sE_t(x)|$ when $x<2R$ and $x\ge 2R$.

Let us consider $x<2R$. Then we notice
$$
\bigg| \Phi'\bigg(\frac{x-sy}{t^{1/(1+\alpha)}}\bigg)\bigg|\le C_1(\alpha).
$$
Indeed, if $\xi$ denotes the argument of $\Phi'$, then
Lemma \ref{l4} yields $|\Phi'(\xi)| \le \frac{|\xi|^\alpha}{\Gamma(1+\alpha)}\le 1$, when  $|\xi|\le 1$ and $|\Phi'(\xi)| \le \frac {C_1(\alpha)}{|\xi|}\le C_1(\alpha)$ for $|\xi|>1$. Hence
$$
|v^-(x,t)-  m \sE_t(x)|\le C_1(\alpha), \qquad\hbox{for }x\le \max\{2R, t^{1/(1+\alpha)}\}.
$$
When $x>\max\{2R, t^{1/(1+\alpha)}\}$, then for $s\in [0,1]$ and $y\in[0,R]$ we have
$|x-sy|/t^{1/(1+\alpha)}>1$. Hence,
$$
\bigg| \Phi'\bigg(\frac{x-sy}{t^{1/(1+\alpha)}}\bigg)\bigg|\le
t^{1/(1+\alpha)} \frac {C_1(\alpha)}{x-sy}.
$$
As a result, we have
\begin{align*}
|I| & \le  \frac{a_0 C_1(\alpha)}{t^{1/(1+\alpha)}}\int_0^\infty
\int_0^1  \frac y{x-sy}\,ds |g(y)|\, dy
%=\frac{a_0}{t^{1/(1+\alpha)}}\int_0^\infty \ln (x/(x-y)) g(y)\,dy \\ &
\le    
\frac{a_0 C_1(\alpha)}{t^{1/(1+\alpha)}}\int_0^\infty \frac y{x-y} |g(y)|\,dy \\ &
\le  \frac{ 2a_0 C_1(\alpha)}{t^{1/(1+\alpha)}}\int_0^\infty \frac y{x} |g(y)|\,dy ,
\end{align*}
where we also used $1/(x-y) \le 2/x$ for $x> 2R.$
\end{proof}

\end{document}
